In summary, we present a novel hardware system for \'etendue expansion and an algorithmic framework for 4D light-field-supervised computer-generated holography. The hardware system includes a multi-source laser array to create a large-\'etendue coherent backlight for the phase SLM and an amplitude SLM for dynamic Fourier-amplitude modulation. The algorithmic framework includes the joint optimization of time-multiplexed amplitude SLM and phase SLM patterns and a memory-efficient, stochastic light field supervision procedure to create 4D light field holograms. We compare our method with a number of \'etendue expansion baselines and verify in simulation and experimentally that our system achieves the highest-quality light field reconstruction results for large \'etendue settings. 

\paragraph{\bf{Limitations and Future Work}} 
We demonstrate our results on a benchtop display setup but futher efforts are required to miniaturize this system. Currently, our multisource laser array is implemented using bulky fiber splitters and could be miniaturized using nanophotonic phased arrays~\cite{sun2013large}. Folding the propagation distance of holograms using optical waveguides could further remove the need of beam splitters and subsequently shrink the form factor, as demonstrated in \cite{jang2024waveguide, yeom2021voxel, lin2018wedge, lin2020waveguide}. We illustrate potential compact designs in the supplemental material. The frame rate of our system is limited by our amplitude display (240 Hz native frame rate). This translates to a $\sim$13.33 Hz frame rate when operating in color-sequential mode with 6-frames time-multiplexing. The frame rate can be improved by using more advanced LCoS displays with frame rate > 720 Hz \cite{lazarev2017lcos}. Real-time synthesis of light field holograms are necessary for practical holographic displays, but is not currently not supported by our system. Extending recent neural network-based hologram synthesis methods \cite{shi2021neural, Shi2022-ih, yang2022diffeng} to work for 4D light field holograms would be an interesting future direction. Finally, we did not attempt to calibrate a neural network--parameterized wave propagation model of our prototype display system, which has been demonstrated to significantly improve experimentally captured holographic image quality for other types of optical configurations~\cite{peng2020neural,choi2022flex}.

\paragraph{\bf{Conclusion}} The novel hardware design and algorithmic framework presented in this work improves the \'etendue of holographic displays and allows for light field holograms synthesis with improved image quality. These help make holographic displays a more practical technology for augmented and virtual reality applications. 